% Figure for Classical Mechanics §B1.1 (Example: payload sliding out along a rotating boom; after Helliwell & Sahakian Example 1.3). Path computed from r = r0 cosh(wt), phi = wt with r0 = 1, R = 6.
\begin{tikzpicture}[>={Stealth[length=2.2mm]}, font=\small, scale=0.62]
  \draw[->, thin, black!55] (-7.6,0) -- (1.8,0) node[right, font=\scriptsize] {$x$};
  \draw[->, thin, black!55] (0,-0.6) -- (0,5.0) node[above, font=\scriptsize] {$y$};
  % boom at t = 0 and at an intermediate time
  \draw[line width=2.2pt, black!30] (0,0) -- (6,0);
  \node[font=\scriptsize, black!60, anchor=north] at (4.2,-0.1) {boom at $t=0$};
  \draw[line width=2.2pt, black!30] (0,0) -- ({6*0.2675},{6*0.9636});
  \node[font=\scriptsize, black!60, anchor=west] at ({6*0.2675},{6*0.9636}) {boom at $\omega t = 1.3$};
  % release position of the boom
  \draw[line width=2.2pt, black!45] (0,0) -- (-4.726,3.696);
  % payload path
  \draw[figB, very thick] (1.000,0.000) -- (1.000,0.036) -- (1.000,0.072) -- (1.000,0.108) -- (1.000,0.145) -- (1.000,0.181) -- (1.000,0.219) -- (0.999,0.257) -- (0.999,0.295) -- (0.998,0.334) -- (0.997,0.374) -- (0.996,0.415) -- (0.994,0.457) -- (0.992,0.500) -- (0.989,0.544) -- (0.986,0.589) -- (0.982,0.636) -- (0.977,0.683) -- (0.971,0.733) -- (0.964,0.783) -- (0.956,0.835) -- (0.946,0.889) -- (0.935,0.944) -- (0.923,1.001) -- (0.908,1.059) -- (0.892,1.119) -- (0.874,1.180) -- (0.853,1.244) -- (0.830,1.309) -- (0.804,1.375) -- (0.776,1.443) -- (0.745,1.513) -- (0.711,1.584) -- (0.673,1.657) -- (0.632,1.731) -- (0.587,1.807) -- (0.538,1.884) -- (0.484,1.962) -- (0.427,2.041) -- (0.365,2.121) -- (0.298,2.202) -- (0.225,2.283) -- (0.148,2.365) -- (0.065,2.448) -- (-0.024,2.531) -- (-0.118,2.613) -- (-0.219,2.695) -- (-0.327,2.777) -- (-0.441,2.858) -- (-0.562,2.938) -- (-0.690,3.017) -- (-0.825,3.093) -- (-0.968,3.168) -- (-1.119,3.241) -- (-1.278,3.310) -- (-1.445,3.377) -- (-1.620,3.440) -- (-1.804,3.499) -- (-1.997,3.553) -- (-2.199,3.603) -- (-2.409,3.646) -- (-2.629,3.684) -- (-2.858,3.716) -- (-3.097,3.740) -- (-3.344,3.756) -- (-3.602,3.763) -- (-3.869,3.762) -- (-4.145,3.751) -- (-4.431,3.729) -- (-4.726,3.696);
  \fill[figB] (1,0) circle (2.2pt) node[below left, font=\scriptsize] {$r_0$};
  % unit vectors and boom force at the intermediate point
  \fill[figB] (0.527,1.899) circle (2.2pt);
  \draw[->, thick, figG] (0.527,1.899) -- ++({1.1*0.2675},{1.1*0.9636}) node[right, font=\scriptsize] {$\hat r$};
  \draw[->, very thick, figO] (0.527,1.899) -- ++({-2.0*0.9636},{2.0*0.2675});
  \node[font=\scriptsize, figO, anchor=north] at (-0.95,1.95) {$F_\varphi\hat\varphi$};
  % release point: velocity components
  \fill[figB] (-4.726,3.696) circle (2.4pt);
  \draw[->, very thick, figG] (-4.726,3.696) -- ++({0.32*5.916*(-0.7880)},{0.32*5.916*0.6157}) node[above, font=\scriptsize] {$v_r\hat r$};
  \draw[->, very thick, figG] (-4.726,3.696) -- ++({0.32*6*(-0.6157)},{0.32*6*(-0.7880)}) node[right, font=\scriptsize] {$v_\varphi\hat\varphi$};
  \draw[->, very thick, figR] (-4.726,3.696) -- ++({0.32*8.426*(-0.9917)},{0.32*8.426*(-0.1284)}) node[left, font=\scriptsize] {$\vec v_f$};
  \node[font=\scriptsize, black!70, anchor=south] at (-3.3,4.0) {release at $r = R$};
  \node[font=\scriptsize, figB, anchor=west] at (1.3,0.9) {$r = r_0\cosh\omega t$};
\end{tikzpicture}
